\documentclass[11pt,a4paper]{article}
\usepackage[T1]{fontenc}
\usepackage[utf8]{inputenc}
\usepackage{lmodern}
\usepackage{amsmath,amssymb,amsthm,mathtools}
\usepackage{esint,graphicx}
\usepackage[a4paper,margin=22mm,headheight=14pt]{geometry}
\usepackage{microtype}
\usepackage{booktabs,longtable,array,enumitem}
\newcolumntype{L}[1]{>{\raggedright\arraybackslash}p{#1}}
\usepackage[dvipsnames]{xcolor}
\usepackage{fancyhdr}
\usepackage{hyperref}

\newtheorem{theorem}{Theorem}[section]
\newtheorem{lemma}[theorem]{Lemma}
\newtheorem{proposition}[theorem]{Proposition}
\newtheorem{corollary}[theorem]{Corollary}
\theoremstyle{remark}
\newtheorem{remark}[theorem]{Remark}
\newtheorem{definition}[theorem]{Definition}
\newtheorem{claim}[theorem]{Claim}
\setlist{nosep,leftmargin=1.5em}
\setlength{\emergencystretch}{2em}
\pagestyle{fancy}
\fancyhf{}
\fancyhead[L]{\small AIM 702}
\fancyhead[R]{\small Revised 7 November 2026}
\fancyfoot[C]{\thepage}
\renewcommand{\headrulewidth}{0.3pt}
\hypersetup{colorlinks=true,linkcolor=MidnightBlue,urlcolor=MidnightBlue,citecolor=MidnightBlue,pdftitle={AIM 702},pdfauthor={Matthew J. Colbrook; results found by AI}}
\begin{document}
\title{Vekua boundary data in triply connected domains}
\author{Matthew J. Colbrook\thanks{The solutions and partial results in this document were found by AI. The supplied research and the separate repository review are AI work; human peer review and formal verification are not claimed.}\\[4pt]
\small Department of Applied Mathematics and Theoretical Physics\\
\small University of Cambridge, Cambridge CB3 0WA, United Kingdom\\
\small\href{mailto:m.colbrook@damtp.cam.ac.uk}{m.colbrook@damtp.cam.ac.uk}}
\date{7 November 2026}
\maketitle
\noindent\textbf{Catalogue scope.} AIM 702; partial progress. \href{https://github.com/MColbrook/AIM/blob/1e564633494c109edfd3ec46da1366adcff925a5/problems/702-vekua-s-harmonic-boundary-problem-in-connectivity-at-least-three.md}{Canonical target at revision \texttt{1e56463}}.


\section{Statement and argument}
\label{sec:702}

The argument below settles the case of three boundary components in
the stated analytic-boundary problem.  The full catalogue target,
which includes every connectivity $n\ge3$, remains unresolved here
for $n\ge4$.

Let $\Omega$ be a bounded connected planar domain bounded by $n$
pairwise disjoint real-analytic Jordan curves.  Label the outer curve
$\Gamma_0$ and the inner curves $\Gamma_1,\ldots,\Gamma_{n-1}$.
Suppose a nonconstant real harmonic function
$u\in C^2(\Omega)\cap C^1(\overline\Omega)$ satisfies
\begin{equation}
 u=c_j,\qquad \partial_\nu u=d_j\qquad\hbox{on }\Gamma_j,
 \label{eq:702-1}
\end{equation}
where the $c_j,d_j$ are real constants and $\nu$ is the outward
unit normal of $\Omega$.
Every $d_j$ is nonzero: constant Dirichlet data and zero normal
derivative on one analytic boundary arc would make $u$ constant by
harmonic reflection and Cauchy uniqueness.

\begin{lemma}[A holomorphic primitive with bounded valence]
Put $q=u_x-iu_y$.  Then $q$ has exactly $n-2$ zeros in $\Omega$,
counted with multiplicity.  There is a single-valued holomorphic
function $F$ on $\Omega$, continuous up to its boundary, such that
\begin{equation}
 F'=q^2,\qquad
 F(z)=-d_j^2\overline z+C_j\quad(z\in\Gamma_j)
 \label{eq:702-2}
\end{equation}
for constants $C_j\in\mathbb C$.
For every $w\notin F(\partial\Omega)$, the number of solutions of
$F(z)=w$ in $\Omega$, counted with multiplicity, is at most $n-1$.
\end{lemma}

\begin{proof}
Orient every boundary component positively with respect to $\Omega$:
the outer component is counterclockwise and the inner components are
clockwise.  Let $\tau=dz/ds$ be the complex unit tangent in this
orientation.  Then $\nu=-i\tau$.
Because $u$ is constant along a boundary component, its tangential
derivative is zero and its gradient is $d_j\nu$.  Consequently
\begin{equation}
 q=d_j\overline\nu=i d_j\overline\tau
 \qquad\hbox{on }\Gamma_j.
 \label{eq:702-3}
\end{equation}
The winding of $q$ about zero is $-1$ on the outer component and $+1$
on each inner component.  The turning-number theorem for smooth
Jordan curves and the argument principle therefore give
\begin{equation}
 \#\{q=0\}\ \hbox{with multiplicity}
 =-1+(n-1)=n-2.
 \label{eq:702-4}
\end{equation}
These uses of the argument principle are legitimate up to the
boundary: locally straighten an analytic boundary arc conformally;
$u-c_j$ has zero boundary values and extends harmonically by odd
reflection.  Thus $q$ extends holomorphically across each boundary
arc and, by $d_j\ne0$, has no boundary zero.

Squaring \eqref{eq:702-3} and using $dz=\tau ds$ gives
\begin{equation}
 q^2\,dz=-d_j^2\overline\tau\,ds
         =-d_j^2\,d\overline z
 \qquad\hbox{on }\Gamma_j.
 \label{eq:702-5}
\end{equation}
In particular every boundary period of $q^2\,dz$ is zero.
The inner boundary cycles generate the first homology of $\Omega$,
so all its periods vanish.  Integrating $q^2\,dz$ defines a
single-valued holomorphic primitive $F$; integrating \eqref{eq:702-5} along
each component proves \eqref{eq:702-2}.

Let $D_0$ be the bounded Jordan domain enclosed by $\Gamma_0$, and
let $D_j$ be the hole enclosed by $\Gamma_j$ for $j\ge1$.
Define the reflected and scaled Jordan domains
\[
 K_j=\{-d_j^2\overline z+C_j:z\in D_j\}.
\]
They need not be disjoint.  The map in this display reverses
orientation.  Hence $F(\Gamma_0)$ winds clockwise around $K_0$, and
each $F(\Gamma_j)$, $j\ge1$, winds counterclockwise around $K_j$.
For $w$ off all these boundary images, the argument principle gives
the exact formula
\begin{equation}
 N_F(w)=\sum_{j=1}^{n-1}{\bf1}_{K_j}(w)-{\bf1}_{K_0}(w).
 \label{eq:702-6}
\end{equation}
The left side counts zeros of $F-w$ with multiplicity, since $F$ has
no poles.  In particular it is nonnegative and is at most $n-1$.
\end{proof}

\begin{theorem}[Triply connected nonexistence]
No nonconstant harmonic function can satisfy \eqref{eq:702-1} on a bounded
triply connected domain with real-analytic Jordan boundary components.
\end{theorem}

\begin{proof}
For $n=3$, \eqref{eq:702-4} says that $q$ has one zero $a$, of multiplicity
one.  Since $F'=q^2$, the holomorphic map $F$ has local degree three
at $a$.  More explicitly,
\[
 F(z)-F(a)=c(z-a)^3+O((z-a)^4),\qquad c\ne0.
\]
Choose a sufficiently small closed disc about $a$ contained in
$\Omega$.  For every $w$ sufficiently close to $F(a)$, Rouch\'e's
theorem gives three solutions of $F(z)=w$ in this disc, counted with
multiplicity.  Choose such a $w$ outside the finite union
$F(\partial\Omega)$ of Jordan curves; this is possible because that
union has empty interior.
Then $N_F(w)\ge3$, contradicting $N_F(w)\le2$ from \eqref{eq:702-6}.
\end{proof}

\begin{corollary}[A restriction at higher connectivity]
Under \eqref{eq:702-1}, if $q$ has a zero of multiplicity $k\ge1$, then
\begin{equation}
 2k+1\le n-1.
 \label{eq:702-7}
\end{equation}
In particular, a nonconstant solution at connectivity $n\ge3$ would
have at least two distinct interior critical points.  At connectivity
four or five all zeros of $q$ would be simple.
\end{corollary}

\begin{proof}
The local degree of $F$ at such a zero is $2k+1$.  The same nearby
value argument bounds it by the global valence $n-1$.
Since the sum of the multiplicities is $n-2$, \eqref{eq:702-7} excludes a
single distinct critical point and gives the last assertions.
\end{proof}

The underlying harmonic problem and the annular classification appear
in \cite[Section 5, Theorem 5.5]{KhavinsonSolyninVassilev2005}.
No sign restriction on the separate fluxes was imposed above: their
squares enter \eqref{eq:702-2}, while their nonzero signs change the argument
of $q$ only by a constant on each boundary component.


\begingroup\small
\begin{thebibliography}{99}
\bibitem{KhavinsonSolyninVassilev2005} D. Khavinson, A. Yu. Solynin and D. Vassilev, \emph{Overdetermined boundary value problems, quadrature domains and applications}, Comput. Methods Funct. Theory \textbf{5} (2005), 19--48, \S5. \url{https://www.unm.edu/~vassilev/quadr.pdf}.
\end{thebibliography}
\endgroup
\end{document}
